%%%%%%%%%%%%%%%%%%%%%%%%%%%%%%%%%%%%%%%%%%%%%%%%%%%%%%%%%%%%%%%%%%%%%%%%
\section{What is proved for finite meanders}\label{sec:proved}
%%%%%%%%%%%%%%%%%%%%%%%%%%%%%%%%%%%%%%%%%%%%%%%%%%%%%%%%%%%%%%%%%%%%%%%%

For finite meanders the proved statements trap the leftover factor $u_n=M_n/\mu^n$ in a window.
The proofs are in Paper~C~\cite{PaperC}; here we explain why they work.

\begin{theorem}\label{thm:sqrt}
There is a constant $C$ such that for every $n\ge1$
\[ \mu^n e^{-C\sqrt n} \;\le\; M_n \;\le\; \frac{\mu^{n+1}}{8n} . \]
\end{theorem}

The upper bound decays like $n^{-1}$, already a power law but with exponent $1$ instead of
$3.42$. The lower bound decays faster than any power of $n$. So the true $u_n$, believed to be
$A\,n^{-3.4201}$, is proved to lie between $e^{-C\sqrt n}$ and $\mu/(8n)$. In logarithmic terms,
\[ -C\sqrt n \;\le\; \log u_n \;\le\; -\log n + \log(\mu/8),\]
while the prediction is $-3.4201\log n+\log A$. Both bounds are new as far as we know. The
lower bound improves one obtained earlier in this project, which had $C\sqrt n\log n$ in the
exponent. Both sides have a short explanation.

\subsection{The upper bound}\label{sec:upperA}

Meanders can be glued: two meanders of orders $a$ and $b$ can be placed side by side on the
river and their roads joined into one, so that the two pieces can be recovered from the result
(Paper~C). This gives $M_{a+b}\ge M_aM_b$. Fekete's lemma says that a sequence $(y_k)$ with
$y_{j+k}\ge y_j+y_k$ has $y_k/k\to\sup_k y_k/k$. Applied to $y_k=\log M_k$, it shows that
$\mu=\lim M_n^{1/n}$ exists, equals $\sup_n M_n^{1/n}$, and so $M_n\le\mu^n$ for every $n$.

A more careful gluing inserts a second meander of order $n$, capped by one extra pair of
crossings (one on each side of it, joined by a new arch below the river), into a gap of a first
meander of order $n$. The result has order $2n+1$, and the insertion can be done in at least
$4n-1$ ways, all of which can be told apart afterwards. So $M_{2n+1}\ge(4n-1)M_n^2$ (Paper~C), and
since $8(2n+1)(4n-1)\ge64n^2$, the numbers $g(n)=8n\,M_n$ satisfy $g(2n+1)\ge g(n)^2$. The upper
bound follows in two lines. Since $(2n+1)+1=2(n+1)$, the numbers $f(n)=\log g(n)/(n+1)$ satisfy
$f(2n+1)\ge f(n)$.
Iterating along $n,\,2n+1,\,4n+3,\ldots$ gives an increasing sequence whose limit is $\log\mu$,
because $g(n)^{1/n}\to\mu$. So $f(n)\le\log\mu$, that is $8nM_n\le\mu^{n+1}$. Intuitively, if
two halves can be glued in about $n$ ways, a meander of order $2n$ is about $n$ times more common
than a pair of halves, and that forces a factor $1/n$.

\subsection{The lower bound}

The lower bound needs the opposite inequality: a large meander must not be much more common than
a pair of smaller ones. Cut it in two. Let the meander have order $m$, so that its sphere has
$2m$ squares. By a form of the planar separator theorem (balanced separators of this size are
classical for planar graphs~\cite{LT,Miller}), the sphere has a closed curve, a
\emph{separator}, that crosses the river and the road at about $\sqrt m$ points in total and
splits the sphere into two pieces of comparable size. Each piece can be completed to a meander:
draw a new closed curve just outside the piece, made of bits of river and road, and then join the
remaining loose ends outside it so that the river and the road each close up into one curve. The
original meander can be recovered from the two completed ones plus a small record. The new
ingredient is that the number of possible records is at most $e^{Bs}$ for a constant $B$, where
$s$ is the number of points where the separator crosses the drawing. Intuitively, the new curve
can be traced again in the completed meander one step at a time, with at most three choices per
step, so the piece is easy to find. Earlier completions had no such curve and needed about
$m^{s}$ records, which produced the extra logarithm.

With $s$ of order $\sqrt m$ this gives
\[ M_m \;\le\; e^{B'\sqrt m}\, M_a M_b \]
for some $a$ and $b$ that each lie between $m/10$ and $9m/10$ and add up to $m$ plus at most a
constant times $\sqrt m$. In terms of $L(n)=n\log\mu-\log M_n=-\log u_n$, which satisfies
$L(n)\ge0$ and $L(n)/n\to0$, this says $L(a)+L(b)\le L(m)+B''\sqrt m$, so one of the two pieces, of size $k$, has
$L(k)/k\le L(m)/m+B''/\sqrt m$. Start from a huge size and keep passing to such a piece until the
size first drops below a target $n$. The sizes shrink geometrically, so the added terms
$B''/\sqrt{\text{size}}$ sum to at most a constant times $1/\sqrt n$, and the starting term
$L(m)/m$ is as small as we like. This gives some $k$ between $n/10$ and $n$ with
$L(k)\le B'''\sqrt n$. Finally, $L(a+b)\le L(a)+L(b)$, so splitting any $n$ into such blocks of
geometrically decreasing sizes gives $L(n)\le C\sqrt n$, which is Theorem~\ref{thm:sqrt}.

The square root has a clear source. Some planar maps, such as square grids, really need
separators with about $\sqrt m$ crossings, and the argument uses one separator for every meander.
To do better, the separator would have to be shorter, which can only hold for a fraction of
meanders, and that fraction would have to be estimated. Paper~C states exactly what such an
estimate would give. It also gives a second criterion, in terms of meanders whose arches are
nested only to a bounded depth, and shows why building such meanders by gluing short ones cannot
meet it. These criteria are proved but not yet independently checked.

\subsection{Prime meanders}

Say that a meander has a \emph{cut} at the gap between crossings $2m$ and $2m+1$ if no arch
below the river passes over the gap, the crossings $2m$ and $2m+1$ are joined by a short arch
above the river, and exactly one other arch above the river passes over the gap. This is a thin
neck of the kind in Section~\ref{sec:shape}, at any position, and the meanders with a cut are
exactly the pairs glued side by side as at the start of Section~\ref{sec:upperA}. A meander is
\emph{prime} if it has no cut. Every meander factors uniquely into prime meanders (a close
relative of known decompositions), and on average a meander has a bounded number of cuts. A small new result is that the chance of a cut at a given
position tends to zero as the distance of that position from both ends grows; its proof has not
yet been independently checked. These facts, an exact identity that relates the power law to
necks at the middle (used in Section~\ref{sec:bridge}), and exact data are in Paper~C.
